\ifx\danslelivre\undefined
\documentclass[10pt]{../fichiers-configuration-latex/smfart}
\usepackage{ifthen}
\newboolean{english}
\newboolean{francais}
\setboolean{english}{false}
\setboolean{francais}{true}
\usepackage{palatino,euler}
\input{../fichiers-configuration-latex/paquets}
\selectlanguage{francais}
\input{../fichiers-configuration-latex/lemme-proposition-theoreme}
\input{../fichiers-configuration-latex/numerotation}
\input{../fichiers-configuration-latex/macros}
\input{../fichiers-configuration-latex/ucs_manquants}
\input{../.cv}


\externaldocument{../13-theoreme-de-finitude/finitude}
\externaldocument{../11-produits-orientes/produits-orientes}

\title{XII. Descente cohomologique orientée}

\begin{document}
\begin{center}
XII. Descente cohomologique orientée\\
Fabrice Orgogozo
\end{center}

\version

\else
\setboolean{english}{false}
\setboolean{francais}{true}
\selectlanguage{francais}
\chapter{Descente cohomologique orientée}
\addtocontents{toc}{par Fabrice Orgogozo}
\begin{center}
Fabrice Orgogozo
\end{center}
\fi

\section{Acyclicité orientée des morphismes propres}

\subsection{}
\label{12-notations invcoho}
L'objet de cette section est de démontrer le théorème \ref{12-invariance cohomologique}
ci-dessous, qui généralise l'invariance par éclatement admissible
du voisinage tubulaire défini à l'aide du produit fibré orienté
(cf. \refext{XI}{11-produit orienté} et \refext{XI}{11-remarques}).
Avant d'énoncer le théorème, fixons quelques notations.
On considère un schéma $S$ cohérent, $g:Y → S$, $π: S′ → S$,
et $a:X → S′$ des morphismes de schémas avec $π$
\emph{propre} et $a$ \emph{cohérent}. 
À ces données sont associés :
\begin{itemize}
\item le morphisme composé $b= π ∘ a$ ;
\item le morphisme $g′: Y′ → S′$ déduit de $g$ par
le changement de base $π$, où $Y′=Y ×_S S′$, et le morphisme de projection $ρ:Y′ → Y$ ;
\item les topos $T=Y\po_S X$ et $T'=Y'\po_{S'} X$
tels que définis en \emph{loc. cit.} par produits
orientés des topos étales associés aux schémas considérés,
ainsi enfin que le morphisme $\fl{ρ}:T'→T$ déduit par fonctorialité
des morphismes $ρ$ et $π$ (\refext{XI}{11-fonctorialité}).
\end{itemize}

\[
\xymatrix{
Y \cartesien \ar[d]_g & Y' \ar[d]^{g'} \ar[l]_{ρ} &  & T'=Y'\po_{S'} X
\ar[d]_{\fl{ρ}} \\
S & S' \ar[l]^{π} & X \ar[l]^{a}
\ar@/^1pc/[ll]^b & T=Y\po_S X
}
\]

\subsection{Énoncés}

\begin{théorème2}
\label{12-invariance cohomologique}
Sous les hypothèses du \ref{12-notations invcoho},
le morphisme $\fl{ρ}$ est acyclique pour les faisceaux de torsion :
pour tout entier $n≥1$ et tout objet $𝒦$
de $\dpluscoefs{T}{𝐙/n}$, le morphisme
d'adjonction $𝒦→\R \fl{ρ}_* \fl{ρ}^*𝒦$ est un isomorphisme.
\end{théorème2}

\begin{théorème2}\label{12-invariance cohomologique non abélienne}
Sous les hypothèses du \ref{12-notations invcoho},
le morphisme $\fl{ρ}$ est $1$-acyclique
pour tout champ ind-fini $𝒞$ sur $T$,
le $2$-morphisme d'adjonction $𝒞→\fl{ρ}_* \fl{ρ}^*𝒞$ est une équivalence de catégories.
\end{théorème2}

\subsection{Démonstrations}

\subsubsection{Réductions}
\label{12-réductions démonstration invcoho}
Soit $𝒦$ comme dans l'énoncé \ref{12-invariance cohomologique}.
Notons $𝒦'$ son image inverse $\fl{ρ}^*𝒦$ sur $T'$.
Le topos $T$ est \emph{cohérent}, cf. \refext{XI}{11-cohérent}.
Il a donc assez de points (\sga{4}{VI}{9.0}) et
ceux-ci sont comme décrits en \refext{XI}{11-points},
c'est-à-dire associés à une paire de points (géométriques) $(y,x)$
de $Y$ et $X$ et à une spécialisation $b(x) ⤳ g(y)$.
En calculant la fibre du morphisme $𝒦→\R \fl{ρ}_* \fl{ρ}^*𝒦$ en un tel
point, on vérifie comme en \refext{XI}{11-ch. base}
que l'on peut supposer $Y,X,S$ et $g$ locaux (pour la topologie étale),
de sorte que le topos $T$ l'est également (cf. \refext{XI}{11-topos locaux}, \refext{XI}{11-can9}).
Sous ces hypothèses supplémentaires, il nous faut montrer que le
morphisme d'adjonction (image inverse)
\[
α : \RΓ(T,𝒦) → \RΓ(T',𝒦')
\]
est isomorphisme.
Considérons les projections $p₁',p₁,p₂$ et $p₂'$
telles que ci-dessous :
$$
\xymatrix{
Y' \ar[d]_{ρ} &  T'=Y'\po_{S'} X \ar[l]_{p₁'} \ar[d]_{\fl{ρ}}
\ar[rd]^{p₂'} & \\
Y & T=Y\po_S X \ar[l]_{p₁} \ar[r]^{p₂} & X
}
$$
de sorte qu'on a en particulier l'égalité tautologique
$\RΓ(T',𝒦')=\RΓ(Y',\R{p₁'}_*𝒦')$.

\subsubsection{Cas où le complexe $𝒦$ provient du schéma $X$}
Supposons qu'il existe un complexe $𝒢 ∈ \dpluscoefs{X}{𝐙/n}$
tel que
\[
𝒦 = p₂^* 𝒢.
\]
Par commutativité du triangle de droite ci-dessus,
on a également $𝒦′={p₂'}^*𝒢$.
Le but $\RΓ(Y',R{p₁'}_*𝒦')$ du morphisme d'adjonction $α$
est donc le complexe $RΓ(Y',\R{p₁'}_*{p₂'}^*𝒢)$,
lui-même isomorphe, d'après \refext{XI}{11-ch. base}
à $\RΓ(Y',{g'}^*\R a_*𝒢)$.
Par propreté de $a$ et le théorème de changement de base
propre, le morphisme d'adjonction
\[
\RΓ(Y,g^*\R π_*\R a_* 𝒢 ) → \RΓ(Y,\R ρ_*{g'}^*\R a_*𝒢)
= \RΓ(Y',{g'}^*\R a_*𝒢)
\]
est un isomorphisme.
Les schémas $Y$ et $S$ étant locaux pour la topologie
étale et $g$ étant un morphisme local, on
a successivement $\RΓ(Y,g^*\R π_*\R a_* 𝒢 )=\RΓ(S,\R b_* 𝒢)$
et $\RΓ(S,\R b_* 𝒢)=𝒢_x$, où $x$ est le point
de $X$ évident. Enfin, le topos $T$ et le
morphisme $p₂$ étant locaux également,
on a $𝒢_x = \RΓ(T,𝒦)=𝒦_t$.
Nous avons exprimé la flèche $α$ comme composée
de morphismes d'adjonction qui sont des isomorphismes. CQFD.

\subsubsection{Cas général}
\label{12-cas général démonstration invcoho}
On ne suppose dorénavant plus que $𝒦$ se descend à $X$.
Considérons cependant le complexe $𝒢=σ^*𝒦$, où $σ$
est la section du morphisme $p₂$ définie en \refext{XI}{11-section canonique}.
(Si $𝒦$ se descend, on retrouve le complexe $𝒢$ ci-dessus.)
D'après \refext{XI}{11-can7}, le morphisme d'adjonction $\R {p₂}_*𝒦→σ^*𝒦$
est un isomorphisme, de sorte que la source $𝒦_t=\RΓ(T,𝒦)=\RΓ(X, \R {p₂}_* 𝒦)$
du morphisme d'adjonction $α$ est isomorphe à $\RΓ(X,𝒢)=𝒢_x=\fl{𝒢}_t$
où l'on pose $\fl{𝒢}={p₂}^*𝒢$. En d'autre termes,
la coünité $\fl{𝒢}→𝒦$ de l'adjonction induit un isomorphisme sur les sections globales
du topos \emph{local} $T=Y\po_S X$. D'après ce qui précède (cas où $𝒦$ se descend),
le morphisme d'adjonction $\RΓ(T,\fl{𝒢}) → \RΓ(T′,\fl{𝒢}′)$ est un
isomorphisme, où $\fl{𝒢}'$ désigne le tiré en arrière
par $p₁'$ de $\fl{𝒢}$ sur le topos $T'$.
Il serait donc suffisant de montrer que le morphisme
$c:\R{p₁'}_*\fl{𝒢}' → \R{p₁'}_*𝒦'$ est un isomorphisme
car en appliquant le foncteur $\RΓ(Y', —)$ on obtient
le morphisme $\RΓ(T′,\fl{𝒢}') → \RΓ(T′,\fl{𝒦}')$ déduit
de la coünité. Par propreté de $π$ donc de $Y′$ sur $Y$,
il suffit même de montrer que la restriction de $c$
à la fibre spéciale $Y′_y$ est un isomorphisme. C'est
ce que nous allons vérifier. (On utilise ici le théorème
de changement de base propre pour les faisceaux de torsion.)
Soient donc $y'$ un point (« géométrique ») de $Y'_y$, $s'$ son image dans $S'$
et calculons la fibre en $y'$ du morphisme $c$.
Notons $Y″$ et $S″$ les localisés $Y'_{(y')}$ et $S'_{(s')}$ ainsi
que $X″$ le produit fibré $X ×_{S′} S″$.
La fibre en $y′$ d'une image directe par $p₁′$
s'identifie à la cohomologie du topos $T″= Y″ \po_{S″} X″$
(à valeurs dans le tiré en arrière).
Si $σ″$ désigne la section canonique (\refext{XI}{11-section canonique})
de la seconde projection $p₂″:T″ → X″$, on a un isomorphisme
canonique (\refext{XI}{11-can7})
\[
\RΓ(T″,-) = \RΓ(X″,{σ″}^*-).
\]
Il en résulte qu'il suffit de montrer que la coünité $\fl{𝒢}″ → 𝒦″$
devient un isomorphisme après application du foncteur $σ″$,
où l'on note $𝒦″$ et $\fl{𝒢}″$ les tirés en arrière sur $T″$.
Ceci résulte immédiatement de la transitivité des images inverses
et de l'égalité $σp₂pσ″=pσ″$ où $p$ désigne
le morphisme $T″ → T$. Cette égalité résulte
à son tour du fait que le morphisme $Y″ → Y$ est
\emph{local}.

\subsubsection{Cas non abélien}
\label{12-remarque-cas-non-abélien}
La démonstration est identique. Pour le théorème de changement de base
propre non abélien, on fait appel à \cite[VII.2.2.2]{CNA@Giraud}.
Signalons que, comme signalé en \refext{XI}{11-remarques} (3)
nous n'avons pas connaissance d'une référence publiée
permettant de justifier le passage à la limite nécessaire au
calcul des fibres.

\section{Descente cohomologique orientée}

\subsection{Topologie orientée des altérations}
\subsubsection{}
Soient $S$ un schéma nœthérien et $ℬ$ la catégorie des diagrammes de schémas
nœthériens $X\dessusdessous{f}{→} S \dessusdessous{g}{←} Y$ et des
morphismes rendant les carrés commutatifs.
Considérons le pseudo-foncteur de $ℬ$ vers la $2$-catégorie
des topos, qui envoie l'objet précédent sur le produit
fibré orienté $Y\po_S X$, où l'on note abusivement $Y$ pour $Y_\et$, etc.
Remarquons que les limites finies (resp. les coproduits) existent
(resp. existent et sont disjoints, universels) dans la
catégorie $ℬ$ ; ils se calculent « terme à terme ». Par exemple,
le produit fibré de $X₁→S₁←Y₁$ et $X₂→S₂←Y₂$ au-dessus de $X→S←Y$
est $(X₁×_X X₂)→ (S₁×_S S₂) ← (Y₁×_Y Y₂)$.

\subsubsection{}
\label{12-topologie alt-orientée}
On considère la topologie orientée des altérations sur $ℬ$ engendrée
par les familles ci-dessous :

\begin{enumerate}
\item $\big((X→S←Y_i) → (X→S←Y)\big)_{i∈I}$, où $(Y_i→Y)_{i∈I}$ est une famille
$\alt$-couvrante ;
\item $\big((X_i→S←Y) → (X→S←Y)\big)_{i∈I}$, où $(X_i→X)_{i∈I}$ est
une famille $\alt$-couvrante ;
\item $\big((X×_S S_i → S_i ← Y×_S S_i) → (X→S←Y)\big)_{i∈I}$, où
$(S_i→S)_{i∈I}$ est une famille $\alt$-couvrante ;
\item $\big((X → S' ← Y×_S S') → (X→S←Y)\big)$, où $S'→S$ est un morphisme propre ;
\item $\big((X×_S S' → S' ← Y)→ (X→S←Y)\big)$, où $S'→S$ est un morphisme étale.
\end{enumerate}

Remarquons que les propriétés des familles
de morphismes (i-v) sont stables par changement de base dans
la catégorie $ℬ$.

\subsection{Énoncés}

\begin{théorème2}\label{12-descente cohomologique}
Soit $(X_{∙}→S_{∙}←Y_{∙})→ (X→S←Y)$ un hyperrecouvrement
pour la topologie orientée des altérations, c'est-à-dire un objet simplicial de $ℬ$
au-dessus de $(X→S←Y)$ tel que pour tout entier $i≥0$, le morphisme
canonique $(X_{i+1}→S_{i+1}←Y_{i+1})→\big(\cosq_i(X_{∙}→S_{∙}←Y_{∙})\big)_{i+1}$
de la catégorie $ℬ$ soit couvrant pour la topologie orientée des
altérations. Sous cette hypothèse, et si l'on note
$\fl{ε}:T_∙=Y_{∙}\po_{S_{∙}} X_{∙} → T=Y\po_S X$ le morphisme d'augmentation, on a les résultats
de descente suivants :
\begin{enumerate}
\item pour tout complexe de torsion borné inférieurement $𝒦$ sur $Y\po_S X$, le morphisme
d'adjonction $𝒦→\R \fl{ε}_*\fl{ε}^*𝒦$ est un isomorphisme.
\item pour tout champ $𝒞$ ind-fini sur $Y\po_S X$, le morphisme
d'adjonction $𝒞→\fl{ε}_*\fl{ε}^*𝒞$ est une équivalence.
\end{enumerate}
\end{théorème2}

\subsubsection{}En d'autres termes, le morphisme $\fl{ε}$ est de descente
cohomologique pour les faisceaux de torsion et les champs ind-finis.
Pour le sens à donner à l'énoncé (ii), on renvoie à \cite{CNA@Giraud} (spécialement
chap. VII, §2.2) et \cite{Alterations@Orgogozo}, §2.

\subsubsection{}
La principale application que nous ferons du théorème précédent
est la formule de changement de base suivante.

\begin{théorème2}\label{12-cb orienté}
Soient $X\dessusdessous{f}{→ }S\dessusdessous{g}{←}Y$ un diagramme de
schémas nœthériens et $ε :S_{∙}→ S$ un hyperrecouvrement pour la topologie des altérations.
Notons $f_{∙}:X_{∙}=X×_{S} S_{∙} →  S_{∙}$
et $g_{∙}:Y_{∙}=X×_{S} S_{∙}→S_{∙}$
(resp. $ε_Y:Y_{∙}=Y×_{S} S_{∙} → Y$)
les hyperrecouvrements pour la topologie des altérations
(resp. pour la topologie orientée des altérations) qui s'en déduisent.
\begin{enumerate}
\item Pour tout complexe de torsion borné inférieurement
$𝒦$ sur $X$, le morphisme
$$
g^*\R f_*𝒦 →  \R{ε_Y}_*\big( {g_{∙}}^*\R{f_{∙}}_{*}
𝒦_{|X_{∙}}\big)
$$
est un isomorphisme.
\item Pour tout champ ind-fini $𝒞$ sur $X$, le morphisme
$$
g^*f_*𝒞 →  {ε_Y}_*\big( {g_{∙}}^*{f_{∙}}_{*}
𝒞_{|X_{∙}}\big)
$$
est un isomorphisme.
\end{enumerate}
\end{théorème2}

\subsection{Démonstration du théorème \ref{12-descente cohomologique}}

Nous commençons par démontrer le (i). Pour la variante
non abélienne (ii), cf. \ref{12-variante non abélienne descente}.

\subsubsection{Réductions}
D'après la théorie générale de la descente — cf. p. ex. \cite[5.3.5]{Hodge3@Deligne} —,
il suffit de démontrer le théorème pour les hyperrecouvrements correspondant à un $0$-cosquelette
associé à une des familles de morphismes du type (i) à (v) (\ref{12-topologie alt-orientée}).
De plus, et pour la même raison, on peut
se contenter du cas particulier où les familles $\alt$-couvrantes
apparaissant dans la description de recouvrement sont,
soit un recouvrement ouvert de Zariski, soit un morphisme
propre et surjectif. Il suffit donc de démontrer les cinq propositions suivantes,
où l'on fixe un objet $(X\dessusdessous{f}{→}S\dessusdessous{g}{←}Y)$
de $ℬ$. Rappelons que $f$ est automatiquement cohérent.
Dans les paragraphes ci-dessous, on reprend les notations
de l'énoncé. En particulier, on note $T$ le topos
$Y\po_S X$, on fixe un entier $n ≥ 1$, un complexe $𝒦 ∈ \Ob  \dpluscoefs{T}{𝐙/n}$
et on note $\fl{ε}$ le morphisme du topos simplicial $T_∙$
vers $T$. On veut montrer que la flèche d'adjonction
\[
α: 𝒦 → \R \fl{ε}_* \fl{ε}^* 𝒦 = \R \fl{ε}_* 𝒦_∙
\]
est un isomorphisme.

\subsubsection{Famille de type (i) : $(X→S←Y′) → (X→S←Y)$}
Soit $Y'→Y$ un morphisme couvrant pour la topologie de Zariski
ou bien propre et surjectif, et $Y_{∙}=\cosq_0^Y (Y')$ son cosquelette.
Par passage aux fibres, on peut supposer les schémas $X$, $Y$ et $S$
ainsi que morphisme $g : Y → S$ locaux. Comme on l'a déjà
dit, le topos $T$ est alors un topos local.
Si $Y'→Y$ est Zariski couvrant, il possède alors une section.
Il en est donc de même du morphisme $T′=Y′ \po_S X → T$, auquel cas
le résultat est connu (cf. \sga{4}{V}{bis 3.3.1}).
Supposons donc le morphisme $Y'→Y$ propre et surjectif.
La flèche $α$ s'identifie au morphisme
\[
𝒦_t=\RΓ(T,𝒦)→\R Γ(T_{∙},𝒦_{∙})=\R Γ(Y_{∙},\R {p_{1 ∙}}_* 𝒦_∙).
\]
Reprenons les notations de \ref{12-réductions démonstration invcoho}.
Par descente cohomologique classique, l'adjonction
$𝒦_t = \RΓ(Y,\R{p₁}_* 𝒦) → \RΓ(Y_∙,(\R{p₁}_* 𝒦)_∙))$ est un isomorphisme.
Par propreté de $ρ:Y′ → Y$, il nous suffit de montrer
que la \emph{restriction au-dessus de $y$} de l'adjonction
\[β:{(\R{p₁}_* 𝒦)}_∙ → {\R p_{1 ∙}}_* 𝒦_∙\]
est un isomorphisme. En effet, le théorème de changement de base
propre nous dit que la source et le but de $α$ se calculent
par restriction aux fibres spéciales.
Que $β_y$ soit un isomorphisme résulte immédiatement du lemme ci-dessous (i),
appliqué aux schémas $Z=Y_i$.

\begin{lemme2}
\begin{enumerate}
Soient $X → S$, $Y → S$ et $Z → S$  des morphismes cohérents.
\item Soit $γ:Z → Y$ un morphisme cohérent.
Notons $\fl{γ}:Z\po_S X→ Y\po_S X$ le morphisme induit.
Le morphisme d'adjonction
\[
γ^*\big(\R{p₁^Y}_*𝒦\big)→ \R{p₁^Z}_*\fl{γ}^*𝒦
\]
est un isomorphisme en tout point de $Z_y$.
\item Soit $δ:Z→X$ un morphisme cohérent.
Notons $\fl{δ}:Y\po_S Z → Y\po_S X$ le morphisme induit.
Le morphisme d'adjonction
\[
δ^*\big(\R{p₂^X}_*𝒦\big)→ \R{p₂^Z}_*\fl{δ}^*𝒦
\]
est un isomorphisme en tout point de $Z$.
\end{enumerate}
\end{lemme2}

\begin{démo}
(i) Si $γ$ est un morphisme local de schéma locaux,
le morphisme $\fl{γ}$ est un morphisme local de topos locaux.
La conclusion en résulte par passage aux fibres en observant
que $Z_{(z)} → Y$ est local si $z$ est un point (géométrique)
localisé sur $Z_y$.
(ii) Même argument : pour tout point (géométrique) $z$
de $Z$, d'image $x$ par $δ$, le morphisme $Y \po_S Z_{(z)} →
Y \po_S X_{(x)}$ est un morphisme local de topos locaux.
\end{démo}

\subsubsection{Famille de type (ii) : $(X′→S←Y) → (X→S←Y)$}
\label{12-X propre}
Quitte à remplacer $p₁$ par $p₂$ et utiliser
le (ii) du lemme ci-dessus, la même démonstration s'applique.
Notons le théorème de changement de base propre
n'apparaît pas directement ici mais est néanmoins
utilisé dans la démonstration du théorème de descente classique
pour $X_{∙}→X$.

\subsubsection{Famille de type (iii) : $(X′=X×_S S′ → S′ ←Y′=Y×_S S′) → (X→S←Y)$}
Notons $S_{∙}=\cosq_0^S (S')$ le cosquelette du morphisme $S′ → S$
et $X_{∙}$, $Y_{∙}$ les schémas simpliciaux qui s'en déduisent.
Par localisation, on se ramène au cas où $S$ est local de sorte que
le cas de la topologie de Zariski est un corollaire de l'existence
d'une section au morphisme $S'→S$. Supposons donc $S'→S$ propre et
surjectif et factorisons $\fl{ε}$ en
\[
Y_{∙}\po_{S_{∙}} X_{∙} \dessusdessous{φ}{→ } Y\po_S X_{∙}
\dessusdessous{π}{→ } Y\po_S X.
\]
On a vu en \ref{12-X propre} que le morphisme $𝒦→\R π_*π^*𝒦$
est un isomorphisme. D'autre part, il résulte du théorème \ref{12-invariance cohomologique} et du fait
que les images directes se calculent cran par cran que 
pour tout complexe $𝒢_{∙} ∈ \dpluscoefs{Y\po_S X_{∙}}{𝐙/n}$ sur $Y\po_S X_{∙}$ — par exemple $π^*𝒦$ —
le morphisme d'adjonction $𝒢_{∙}→\R φ_*φ^*𝒢_{∙}$ est un isomorphisme.
Le résultat en découle.

\subsubsection{Famille de type (iv) : $(X → S' ← Y′= Y×_S S') → (X→S←Y)$}
On suppose ici que le morphisme $X → S$ se factorise
à travers un morphisme \emph{propre} $S′ → S$, et on définit $S_∙$ et $Y_∙$ comme ci-dessus.
Remarquons que pour chaque entier $i ≥ 0$, le topos $Y_i\po_{S_i} X$ est
le topos associé au produit fibré itéré dans la catégorie $ℬ$ du
morphisme $(X→S'←Y′)→ (X→S←Y)$.
Factorisons le morphisme $\fl{ε} : Y_{∙}\po_{S_{∙}} X → Y\po_S X$
en
\[
Y_{∙}\po_{S_{∙}} X \dessusdessous{φ}{→ } (Y\po_S X)_{\mathrm{const}} \dessusdessous{π}{→ } Y\po_S X,
\]
où $(Y\po_S X)_{\mathrm{const}}$ est le topos simplicial constant.
Il résulte du théorème \ref{12-invariance cohomologique} appliqué pour
chaque $i$ aux morphismes $a:X\dessusdessous{\text{diag}}{→}S_i$
et $b:S_i→S$ que l'adjonction $\Id→\R φ_*φ^*$ est un isomorphisme. On s'est donc ramené
au cas trivial où $S'=S$.

\subsubsection{Famille de type (v) : $(X′= X×_S S' → S' ← Y)→ (X→S←Y)$}
On suppose ici que le morphisme $Y → S$ se factorise
à travers un morphisme \emph{étale} $S′ → S$.
Notons à nouveau $S_{∙}$ le $0$-cosquelette de ce morphisme
et $X_{∙}=X ×_S S_{∙}$ le schéma simplicial qui s'en déduit.
Comme ci-dessus, notons que pour chaque entier $i ≥ 0$ le topos $Y\po_{S_i} X_i$
est le topos associé au produit fibré itéré dans la catégorie $ℬ$ du
morphisme $(X′ →S'←Y)→ (X→S←Y)$. Pour montrer que la flèche
d'adjonction $α: 𝒦 → \R \fl{ε}_* \fl{ε}^* 𝒦$ est un isomorphisme, on peut supposer — par passage aux
fibres — que le schéma $S$ est strictement hensélien de
sorte que le morphisme $S'→S$, et par conséquent $(X×_S S' → S' ← Y)→
(X→S←Y)$, aient une section. Le résultat est alors évident.

\subsubsection{Variante non abélienne}
\label{12-variante non abélienne descente}
Modulo la difficulté signalée ci-dessus (\ref{12-remarque-cas-non-abélien}), il suffit de remplacer les références
\cite[5.3.5]{Hodge3@Deligne} et \sga{4}{V}{bis 3.3.1}
par \cite[2.5,2.8]{Alterations@Orgogozo}. 

\subsection{Démonstration du théorème \ref{12-cb orienté}}

\subsubsection{Cas des complexes}
Soit $𝒦$ un complexe sur $X$ comme dans l'énoncé,
et notons $\fl{𝒦}$ son image inverse par la seconde
projection $p₂:Y\po_S X → X$.
Il résulte du théorème \ref{12-descente cohomologique} (i)
que l'adjonction $α: 𝒦 → \R \fl{ε}_* \fl{ε}^*𝒦$
est un isomorphisme. (On rappelle que $\fl{ε}$
désigne le morphisme d'augmentation $Y_{∙}\po_{S_{∙}} X_{∙} → Y\po_S
X$.)

D'autre part, on a la chaîne d'isomorphismes :
\begin{align*}
g^*\R f_*𝒦					& ⥲ \R{p₁}_*\fl{𝒦} \\
\R{p₁}_*α : \R{p₁}_*\fl{𝒦}			& ⥲ \R{p₁}_*{\R \fl{ε}}_* {\fl{ε}}^*\fl{𝒦} \\
\R{p₁}_*\R{\fl{ε}}_* {\fl{ε}}^*\fl{𝒦}	& = \R {ε_Y}_* \R {p_{1∙}}_* \fl{𝒦}_{∙}  \text{, où } \fl{𝒦}_{∙}={\fl{ε}}^*\fl{𝒦} \\
\R {ε_Y}_* \R {p_{1∙}}_* \fl{𝒦}_{∙}		& ⭉ \R{ε_Y}_* g_{∙}^* \R {f_{∙}}_* \fl{𝒦}_{∙}
\end{align*}
Les premier et dernier isomorphismes résultent de \refext{XI}{11-ch. base}
et le troisième de la fonctorialité des images directes.

\subsubsection{Cas des champs}
Même démonstration. 

\ifx\danslelivre\undefined
\bibliography{../fichiers-configuration-latex/bibliographie}
\bibliographystyle{../fichiers-configuration-latex/nomcomplet-date}
\end{document}
\fi

